%! AP Statistics — Unit 3, Lesson 3.7 — Student Guided Notes
\documentclass[10pt]{article}
\usepackage{apstats-article}
\usepackage{apstats-boxes}

\begin{document}

\pageheader{Unit 3, Lesson 3.7}{Guided Notes}
\namedateperiod

\vspace{0.12in}

\begin{objectivebox}
\small By the end of this lesson, I will be able to\dots\\[4pt]
\writeline\\[1pt]
\writeline\\[1pt]
\writeline
\end{objectivebox}

\vspace{0.1in}

\begin{vocabbox}
\termblanklong{Statistical inference}
\termblanklong{Statistically significant}
\termblanklong{Causal conclusion}
\termblanklong{Generalization}
\termblanklong{Scope of inference}
\end{vocabbox}

\vspace{0.1in}

\begin{notesbox}{Interpreting Experiment Results (VAR-3.E)}
\small
\textbf{Statistical inference:} attributing conclusions from data to the
\blank{6cm} from which the data were collected.

\vspace{8pt}
\textbf{Statistical significance:} an observed difference is \textbf{statistically
significant} if it is too \blank{3cm} to be explained by \blank{3cm} alone
(VAR-3.E.2).

\vspace{8pt}
\textbf{Causal conclusion} (VAR-3.E.3): If results are statistically significant in
an experiment with random assignment, we can conclude the \blank{5cm} \emph{caused}
the effect.
\end{notesbox}

\vspace{0.1in}

\begin{notesbox}{Scope of Inference Table}
\small
\renewcommand{\arraystretch}{2.2}
\begin{tabularx}{\linewidth}{@{} c | X | X @{}}
 & \textbf{Random Selection used} & \textbf{No Random Selection} \\
\hline
\textbf{Random Assignment used} &
Can conclude \blank{3cm} \textbf{AND} can \blank{3cm} to the population. &
Can conclude \blank{3cm}, but can only generalize to \blank{3cm}. \\
\hline
\textbf{No Random Assignment} &
Cannot conclude \blank{3cm}, but can \blank{3cm} to the population. &
Cannot conclude \blank{3cm} and cannot \blank{3cm} beyond the sample. \\
\end{tabularx}

\vspace{8pt}
\textbf{Key rules:}
\begin{itemize}[itemsep=3pt]
  \item Random \textbf{assignment} $\Rightarrow$ \blank{4cm}
  \item Random \textbf{selection} $\Rightarrow$ \blank{4cm}
\end{itemize}
\end{notesbox}

\vspace{0.1in}

\begin{practicebox}
\small For each study, complete the scope-of-inference analysis.

\renewcommand{\arraystretch}{1.9}
\begin{tabularx}{\linewidth}{@{} X c c p{2.4cm} p{2.8cm} @{}}
\textbf{Study} & \textbf{Rand. assign.?} & \textbf{Rand. select.?}
  & \textbf{Causation?} & \textbf{Generalize to?} \\
\hline
500 randomly selected adults surveyed; found association between exercise and mood.
& Yes / No & Yes / No & Yes / No & \blank{2.2cm} \\
120 volunteers randomly assigned to meditation or control; significant reduction in
stress found.
& Yes / No & Yes / No & Yes / No & \blank{2.2cm} \\
200 randomly selected patients randomly assigned to new drug or placebo; significant
improvement found.
& Yes / No & Yes / No & Yes / No & \blank{2.2cm} \\
\end{tabularx}
\end{practicebox}

\end{document}
